\section{Analytical preparation must preserve the population and its information clock}
\label{exec19:duckdb-preparation}

A servicing analyst may ask whether an additional contact improves recovery,
then construct the analysis table by joining account records, contact events
and repayment observations. The resulting query can execute successfully while
changing the population, discarding unobserved outcomes or incorporating
information that did not exist when contact was chosen. The database is
therefore part of the causal workflow's measurement boundary. Its useful
product role is to make a declared preparation repeatable and inspectable;
query success cannot establish the causal premises of the subsequent analysis.

DuckDB is a concrete candidate for local analytical preparation. The selected
historical source is version 1.4.0, released on 16 September 2025, at commit
\texttt{b8a06e4a2267}. This choice fixes a reproducible object for discussion;
it is not a claim about the newest release, present support or any installed
binary. The inspected C++ interface distinguishes a materialized
\texttt{Query} call from streaming and pending paths. Result objects expose
column information, errors and potentially chained results, while fetching a
materialized result consumes its contents. These interfaces suggest a narrow
first integration: a typed request produces one bounded, completed table and
an explicit completion status. They do not justify treating every query handle,
first batch or returned row as a completed analysis. The Python client is a
separate wrapper boundary whose conversions and installed dependencies require
their own inspection.\citep{exec19_duckdb_source}

\subsection{A missing outcome is not an observed failure}
\label{exec19:duckdb-missingness}

Consider an authored six-account example with three contacted and three
uncontacted accounts. The contacted outcomes are $(1,\mathrm{missing},
\mathrm{missing})$ and the uncontacted outcomes are $(1,0,0)$, where one means
recovery. These are stipulated binary observations, without a randomization,
exchangeability or missingness assumption. The contrast among observed
outcomes is
\begin{equation}
 \widehat d_{\mathrm{observed}}=1-\frac13=\frac23.
 \label{exec19:duckdb-observed}
\end{equation}
If the two missing contacted outcomes are denoted $u,v\in\{0,1\}$, the
full six-account descriptive contrast instead satisfies
\begin{equation}
 d(u,v)=\frac{1+u+v}{3}-\frac13
       =\frac{u+v}{3}\in\left\{0,\frac13,\frac23\right\}.
 \label{exec19:duckdb-completions}
\end{equation}
Both endpoints are attainable, so $[0,2/3]$ is the sharp endpoint range for
this stipulated finite completion problem. The feasible set has three values;
it is not every real number in that interval. Neither the complete-case
contrast nor the completion set is a causal effect without additional
assignment and observation assumptions. Exact rational enumeration reproduces
all four binary completions.\citep{exec19_duckdb_reference}

The economic consequence of the transformation is visible before any bank
experiment. Suppose solely for illustration that a unit change in this
descriptive contrast is valued at USD 300 and an additional contact policy
costs USD 40 per case. The resulting arithmetic is
\begin{equation}
 V(u,v)=300d(u,v)-40\in\{-40,60,160\}\quad\text{USD per case}.
 \label{exec19:duckdb-value}
\end{equation}
Using only observed outcomes produces the optimistic plug-in value USD 160;
the declared completion problem also permits minus USD 40. This is an
illustration of how an unacknowledged denominator choice can reverse an
economic recommendation, not evidence of treatment value or bank ROI. A
missingness model, additional observations or a deliberately restricted target
population would change the problem and must be stated before interpretation.
\citep{exec19_duckdb_reference}

The selected upstream aggregate tests reinforce why the interface must retain
these distinctions. On an empty input, the tests expect zero counts but NULL
for sums, averages and several other aggregates; grouped all-null values also
produce NULL aggregates. Those are inspected upstream expectations, not native
test results obtained here. A zero-row relation, two rows with no observed
outcome, and two observed zeros therefore require different records. The
proposed record carries total rows, observed outcomes and nullable summaries;
it never silently replaces a missing mean with zero. NaN, overflow and invalid
type conversions need additional explicit cases rather than being folded into
the SQL NULL convention.\citep{exec19_duckdb_source,exec19_duckdb_reference}

\subsection{Joins and clocks define the observation being analyzed}
\label{exec19:duckdb-keys}

For a second exact example, let two contacted accounts have observed amounts
USD 100 and USD 0, and two uncontacted accounts both have USD 0. The
account-level descriptive difference is USD 50. Suppose a one-to-many contact
join duplicates the positive contacted account three times while leaving each
other account once. The joined table has six rows and a contacted mean of
USD 75. No account outcome changed: the join changed the weights. The input
and output keys, expected join cardinality and analytical unit must therefore
be part of the preparation request. A deliberate contact-event target is
possible, but it is a different target and does not turn six joined rows into
six independent accounts.\citep{exec19_duckdb_reference}

Time ordering creates a separate problem. In the authored reference, an
account has a feature with value zero on day one, an action on day two and a
feature with value one on day three. Selecting the latest unrestricted record
returns the postaction value; selecting the latest strictly preaction record
returns zero. A production definition must additionally distinguish event time
from availability time: a backdated event received after the decision was not
necessarily available to the decision maker. The proposed table carries both
clocks, the action cutoff, time-zone interpretation and an explicit tie-break
for events with equal timestamps. This clock discipline is an authored
requirement; the small reference does not validate a bank's record history.

Row position is an unsuitable replacement for those keys. The selected tests
make null ordering configurable, and an insertion-order test explicitly
disclaims which tuples an unordered limit selects when insertion-order
preservation is disabled. A reproducible export consequently specifies a total
order with unique tie-breaks and a declared NULL convention. All 24 permutations
of the four-account reference return the same rows after sorting by its unique
account key. That exact property belongs to the reference construction; native
parallel execution, collations and timestamp conversions remain future
conformance cases.\citep{exec19_duckdb_source,exec19_duckdb_reference}

\subsection{Local execution needs explicit permissions and completion}
\label{exec19:duckdb-controls}

The selected source initially enables external access. When external access is
disabled on a running instance, its setting handler adds already attached
database and write-ahead-log paths to allowed paths and permits a configured
temporary directory. Selected storage and attachment tests continue to permit
writes or checkpoints to attached databases after this setting is disabled.
It is therefore neither a read-only guarantee nor an operating-system sandbox.
The proposed worker needs independently defined filesystem permissions,
reviewed attachments and scratch storage, even when that database option is
false. Its input copies must be immutable under the declared execution policy.
\citep{exec19_duckdb_source}

Extension behavior is also a build and runtime question. The selected CMake
options default automatic installation and loading to false, but compile-time
definitions can change the configuration defaults. The base extension
configuration includes core functions and Parquet, and conditionally includes
jemalloc under its platform rules. Inspecting those declarations does not prove
the configuration of a wheel, executable or embedded library. For the initial
proposal profile, automatic extension installation and loading should be
disabled and any required artifact explicitly selected and bound. An agent
may propose a typed preparation request; it should not acquire unrestricted SQL
and installation authority as a side effect. The inspected autoload helper
catches exceptions and can return false. Such a dependency failure is not
evidence that an intervention is unidentified or that an estimated effect is
zero.\citep{exec19_duckdb_source}

Transactions do not remove the need for those boundaries. The C++ begin,
commit and rollback wrappers inspect errors and throw them; a selected upstream
test expects a conflict for overlapping updates. These observations support
explicit transaction-error handling but do not prove a complete isolation
regime or an immutable snapshot of files outside the database. The preparation
manifest binds input bytes separately from the transaction identifier and
records whether an output is complete. A timeout, cancellation, incomplete
stream or later error invalidates a claim of completed output. Materialization
simplifies the initial completion boundary while increasing memory exposure;
row and byte limits, disk use and cancellation costs must be measured on the
selected prototype rather than assumed from the interface name.
\citep{exec19_duckdb_source}

\subsection{Reuse the original engine under a bounded proposal contract}
\label{exec19:duckdb-route}

The preferred proposal route retains the original engine behind that typed
preparation boundary. A port would reproduce SQL, optimization, storage and
conversion obligations without a demonstrated platform need. The selected
historical build declares C++11, twelve direct non-amalgamation third-party
subdirectories, additional conditional test directories and extension inputs.
The root MIT notice is not the complete distribution inventory. The inspected
RE2 and Zstd notices are BSD-style, utf8proc carries MIT and Unicode data
notices, and MbedTLS declares Apache-2.0 or GPL-2.0-or-later alternatives. These
four inspected third-party notices illustrate scope differences; they do not
clear every dependency, wrapper, extension or actual distribution arrangement.
\citep{exec19_duckdb_source}

A transparent explicit-file implementation remains the small-case reference.
SQLite is a separate SQL/storage candidate; dataframe and columnar libraries
are separate preparation candidates. Each alternative requires its own type,
join, missingness, ordering and resource review. The choice among them should
follow the declared workload and measured end-to-end costs, including frozen
inputs, conversions, failed requests, maintenance and retained evidence. No
unmeasured speed, feature parity or maintenance advantage is assumed here.

The proposal choice is consequently concrete while executable admission is
deferred. Before admitting the selected engine/client path, a prototype must
bind its complete distribution and notices, exercise native empty/all-missing/
zero cases, duplicate joins, information cutoffs, total ordering and precision,
and compare errors, cancellations, external-file changes and extension
behavior with the typed contract. The analyst owns the scientific population
and causal assumptions; the adapter maintainer owns conversion, execution and
resource behavior; the distribution owner reviews the actual retained software
and notices. None of these public or synthetic prototype steps requires bank
data. Institution-specific assignment, audit truth and economic value remain
distinct later questions. The database prepares the evidence; it does not
certify the causal meaning of that evidence.
